\documentclass[11pt]{article}

% —— Core packages ——
\usepackage{amsmath,amssymb}
\usepackage{stmaryrd}
\usepackage[utf8]{inputenc}
\usepackage{textcomp}
\usepackage{times}
\usepackage[margin=1in]{geometry}
\usepackage{microtype}
\emergencystretch=1em
\usepackage[most]{tcolorbox}
\tcbuselibrary{theorems,skins,breakable}

% —— Modal logic and Greek ——
\DeclareUnicodeCharacter{03B9}{\ensuremath{\iota}}   % ι iota
\DeclareUnicodeCharacter{25A1}{\ensuremath{\Box}}     % □ necessity
\DeclareUnicodeCharacter{25C7}{\ensuremath{\Diamond}} % ◇ possibility
\DeclareUnicodeCharacter{03BB}{\ensuremath{\lambda}}  % λ lambda (precaution)
\DeclareUnicodeCharacter{03C6}{\ensuremath{\varphi}}  % φ phi (precaution)
\DeclareUnicodeCharacter{00AC}{\ensuremath{\neg}}     % ¬ negation
\DeclareUnicodeCharacter{2227}{\ensuremath{\wedge}}   % ∧ (already listed but verify)
\DeclareUnicodeCharacter{22A2}{\ensuremath{\vdash}}   % ⊢ turnstile (precaution)

\newenvironment{definitionbox}[1][]{%
  \begin{center}
  \begin{tcolorbox}[
    enhanced,
    breakable,
    width=0.8\textwidth,
    colback=gray!2,
    colframe=gray!60,
    coltitle=black,
    fonttitle=\bfseries,
    title={Definition: #1},
    boxed title style={
      colback=gray!3,
      colframe=black!60,
      sharp corners,
      boxrule=0.5pt,
    },
    boxrule=0.5pt,
    sharp corners,
    left=6pt,right=6pt,top=6pt,bottom=6pt,
    before skip=12pt,
    after skip=12pt,
  ]}%
  {\end{tcolorbox}\end{center}}

% —— Tables ——
\usepackage{array}
\usepackage{longtable}
\usepackage{booktabs}

\usepackage{calc}
\newcolumntype{C}[1]{>{\raggedright\arraybackslash}p{\dimexpr#1\linewidth-4\tabcolsep\relax}}

% —— Fix paragraph indent after lists ——
\usepackage{parskip}
\setlength{\parindent}{1.5em}
\setlength{\parskip}{0.5em}

% —— Section formatting ——
\usepackage{titlesec}
\usepackage{needspace}
\usepackage{etoolbox}

\setcounter{secnumdepth}{3}

\titleformat{\section}[hang]
  {\normalfont\Large\bfseries}
  {\thesection.}{1em}{}
\titleformat{\subsection}[hang]
  {\normalfont\large\bfseries}
  {\thesubsection.}{1em}{}
\titleformat{\subsubsection}[hang]
  {\normalfont\normalsize\bfseries}
  {\thesubsubsection.}{1em}{}

\titlespacing*{\section}{0pt}{1.8em}{0.8em}
\titlespacing*{\subsection}{0pt}{1.4em}{0.6em}
\titlespacing*{\subsubsection}{0pt}{1em}{0.4em}

\newcommand{\sectionbreak}{\vspace{0.8\baselineskip}}

% —— TOC formatting ——
\usepackage{tocloft}
\renewcommand{\cftsecfont}{\small}
\renewcommand{\cftsubsecfont}{\small}
\renewcommand{\cftsubsubsecfont}{\small}
\renewcommand{\cftsecpagefont}{\small}
\renewcommand{\cftsubsecpagefont}{\small}
\renewcommand{\cftsubsubsecpagefont}{\small}
\setlength{\cftbeforesecskip}{0.3em}
\setlength{\cftbeforesubsecskip}{0.1em}
\setlength{\cftbeforesubsubsecskip}{0.1em}

% —— Lists ——
\usepackage{enumitem}
\setlist{
  nosep,
  leftmargin=2em,
  topsep=0.4em,
  partopsep=0pt,
  parsep=0.2em,
  itemsep=0.2em
}
\setlist[itemize,1]{label=\textbullet}
\setlist[itemize,2]{label=\textendash}

% —— Figures ——
\usepackage{graphicx}
\usepackage{float}
\setkeys{Gin}{keepaspectratio,width=0.7\linewidth}
\makeatletter
\def\fps@figure{H}
\makeatother
\graphicspath{{figures/}}

% —— URL/DOI wrapping ——
\usepackage{xurl}

% —— CSL bibliography ——
\newlength{\cslhangindent}
\setlength{\cslhangindent}{1.5em}
\newenvironment{CSLReferences}[2]%
  {\begin{enumerate}[label={[\arabic*]},leftmargin=\cslhangindent,itemsep=0pt]%
   \raggedright\sloppy
   \renewcommand{\bibitem}[2][]{\item}}%
  {\end{enumerate}}
\providecommand{\CSLLeftMargin}[1]{\item[#1]}
\providecommand{\CSLRightInline}[1]{#1}

% —— Hyperref (must be near last) ——
\usepackage[colorlinks=true, linkcolor=black, urlcolor=blue, citecolor=blue]{hyperref}
\urlstyle{same}

% —— Fix for Pandoc/longtable "none" counter ——
\makeatletter
\@ifundefined{c@none}{%
  \newcounter{none}%
}{}
\makeatother
\renewcommand{\thenone}{}          % don't actually print a number
\providecommand*{\theHnone}{\arabic{none}} % for hyperref's anchors

% —— Paragraph formatting (block style) ——
\setlength{\parindent}{0pt}
\setlength{\parskip}{0.6em}
\linespread{1.05}
\raggedbottom

% —— Block quotes ——
\renewenvironment{quote}{%
  \par\vskip 0.5em
  \leftskip=2em\rightskip=2em\small\itshape
  \noindent\ignorespaces
}{\par\vskip 0.5em}

% —— Pandoc compatibility ——
\providecommand{\tightlist}{}
\providecommand{\pandocbounded}[1]{#1}

% —— Citeproc fixes ——
\makeatletter
\providecommand{\citeproctext}{} % do NOT force a constant label
\let\@oldprotect\protect
\def\@citeprocprotection{\let\protect\@oldprotect}
\newcommand{\citeproc}[2]{#2}
\protected\def\citeprocfootnote#1{\footnote{#1}}
\makeatother

% —— Special character definitions ——
\DeclareUnicodeCharacter{2203}{\ensuremath{\exists}}
\DeclareUnicodeCharacter{2200}{\ensuremath{\forall}}
\DeclareUnicodeCharacter{2227}{\ensuremath{\wedge}}
\DeclareUnicodeCharacter{2228}{\ensuremath{\vee}}
\DeclareUnicodeCharacter{2192}{\ensuremath{\rightarrow}}
\DeclareUnicodeCharacter{2194}{\ensuremath{\leftrightarrow}}
\DeclareUnicodeCharacter{2234}{\ensuremath{\therefore}}
\DeclareUnicodeCharacter{2248}{\ensuremath{\approx}}
\DeclareUnicodeCharacter{2208}{\ensuremath{\in}}
\DeclareUnicodeCharacter{2260}{\ensuremath{\neq}}
\DeclareUnicodeCharacter{2713}{\checkmark}
\DeclareUnicodeCharacter{2212}{-}

% —— Angle brackets for LPCs ——
\newcommand{\llangle}{\langle\!\langle}
\newcommand{\rrangle}{\rangle\!\rangle}
\DeclareUnicodeCharacter{27EA}{\ensuremath{\llangle}}
\DeclareUnicodeCharacter{27EB}{\ensuremath{\rrangle}}

% —— Custom commands for frequent notation ——
\newcommand{\Rfr}[2]{\ensuremath{\mathcal{R}(#1, #2)}}
\newcommand{\LPC}[1]{\ensuremath{\llangle#1\rrangle}}

% —— Document metadata ——
\title{%
}
\author{Florin Cojocariu} 
\date{}

\begin{document}
\begin{center}
\vspace*{0.5cm}
{\large UNIVERSITATEA DIN BUCUREȘTI}\\[0.3em]
{\large FACULTATEA DE FILOSOFIE}

\vfill

{\large\bfseries Lucrare de disertație}

\vspace{2em}

{\LARGE\bfseries From Local to Portable Rigidity}\\[1em]
{\Large Demonstrative Reference and the Cognitive Ground of Rigid Designation}

\vfill
\end{center}
\noindent
\begin{minipage}[t]{0.45\textwidth}
\raggedright
\textbf{Coordonator științific:}\\
Prof. Dr. Mircea Dumitru
\end{minipage}%
\hfill
\begin{minipage}[t]{0.45\textwidth}
\raggedleft
\textbf{Absolvent:}\\
Florin Cojocariu
\end{minipage}
\vspace{1.5cm}
\begin{center}
{\large București}\\[0.3em]
{\large 2026-05-24}
\end{center}
\clearpage

\section*{Abstract}\label{abstract}
\addcontentsline{toc}{section}{Abstract}

The dissertation argues that Kripke's account of rigid designation rests on a cognitive precondition his semantics presupposes but does not theorize. In Kripke's natural-language gloss of~\(\forall x \forall y\,(x = y \rightarrow \Box(x = y))\)~--- ``for every object~x~and object~y'' --- the word ``object'' spans three distinct things: mathematical variables rigid by assignment, encountered instances rigid within an encounter, and named individuals rigid across absence. The formula's validity holds for the first; the philosophical claim concerns the third. Kripke argues for that claim independently and is aware the formal result concerns variables, not names; what he does not provide is a cognitive or metasemantic account of how ordinary speakers achieve the referential stability his semantics attributes to names. That account is what this dissertation supplies.

The thesis is that rigid name use can be illuminated metasemantically by comparison with \(x^o\)-mode tracking across absence, where \(x^o\) marks the use of a word in coordination with an experienced instance \(\{X\}\) and \(x^c\) marks its use within broader linguistic relations. The proposal distinguishes two levels of linguistic-cognitive operation: at~\(x^o\), a speaker's use is anchored to a particular through live coordination; at~\(x^c\), terms function within broader inferential and testimonial networks, with no experienced instance that it refers to. The dissertation argues that the modal force characteristic of rigid designation depends on the former, while the epistemic work of identification and convergence is largely carried out by the latter.

The argument runs through Russell, Kripke, and the post-Kripkean cognitive literature. Russell identified the mechanism --- demonstrative reference under acquaintance --- and then restricted it to immediate sense-data on Cartesian grounds. Kripke extended the scope of rigidity to natural-language names but inherited Russell's restraint by leaving the mechanism unspecified. Existing cognitive and metasemantic accounts --- including Evans's discriminating conception, Recanati's mental files, and Devitt's causal-historical picture --- each capture part of the explanatory burden. My claim is not that they ignore the problem, but that they do not yet explain rigid designation in terms of a single primitive coordination that scales from demonstrative encounter to socially transmitted name use. The dissertation locates the mechanism Russell discovered and Kripke extended in \(x^o\)-mode tracking, with baptismal lifting carrying that tracking across absence through social practice.

\vfill
\noindent\rule{\textwidth}{0.4pt}\\[2pt]
{\tiny florin.cojocariu@s.unibuc.ro, \today}

\end{document}